\newif\ifglobal

%%%%%%%%%%%%%%%%%%%%%%%
%%% Compile to view %%%
%%%%%%%%%%%%%%%%%%%%%%%

%%%%%%%%%%%%%%%%%%%%%%%%%%%%%%%%%%%%%%%%%%%%%%%%%%%
%%% Readme:                                     %%%
%%% https://www.overleaf.com/read/bwdjvfjksvjy/ %%%
%%%%%%%%%%%%%%%%%%%%%%%%%%%%%%%%%%%%%%%%%%%%%%%%%%%

% >>> M?: marks a module setting number or important notice

% >>> M4: Set {./relative-path} to external files for this module <<<
\def \modulefiles {./30-GLITCH}
% To add an external file, use:
% (Replace '---' with actual filename)
% '\input{\modulefiles/tables/---.tex}' for tables
% '\includegraphics{\modulefiles/figures/---}' for figures
% '\subfile{\modulefiles/---}' for register files


% >>> M1: This module is in global project? <<<
% 'YES' - uncomment; 'NO' - comment out
\globaltrue

\ifglobal
    \documentclass[main\_\_CN.tex]{subfiles}
   
\else
    % Font "FandolSong-Regular" used by ctex does not have GSUB table.
    % It causes warnings that do not affect compilation and can be ignored.
    \PassOptionsToPackage{quiet}{fontspec}

    \documentclass[openany, 10pt]{book}
   
    % >>> M2: In ./00-shared/preamble-trm-module, also find
    %         settings for visibility of todo notes, watermark <<< 
    \usepackage{preamble-trm-module}


    % Enable Chinese typesetting
    \usepackage{ctex}

    % >>> M3: Temporary labels in Chinese
    %         you can add your own labels there <<<
    \myexternaldocument{./00-shared/temp-labels-cn}

\fi %\ifglobal


\setcounter{secnumdepth}{4}


% >>> M5: If this module needs any updates in
% './00-shared/preamble-trm-module.sty'
% (include additional package, etc.), contact your module PM
% or create the file './00-shared/changelog.tex' make the
% necessary updates yourself and document them in this file <<<


\begin{document}


% Sets language into which pre-defined bits of text are translated
\selectlanguage{chinese}

% Do not delete this block
\ifglobal
% Add the following front matter if
% this module is outside of global project
\else
    \InputIfFileExists{readme}{}{}
    {\let\clearpage\relax \listoftodos}
    \tableofcontents
    %\listoftables
    %\listoffigures
\fi


%%%%%%%%%%%%%%%%%%%%%%%%%
%%% TRM Module Begins %%%
%%%%%%%%%%%%%%%%%%%%%%%%%
\hypertarget{glitch}{}
\chapter{时钟毛刺检测}
\label{mod:glitch}


\section{概述}

为提升 ESP32-C3 的安全性能，防止攻击者通过给外部晶振 XTAL\_CLK 附加毛刺，使芯片进入异常状态，从而实施对芯片的攻击，ESP32-C3 搭载了毛刺检测模块(CLK Glitch\_Detect)，用于检测从外部晶振输入的 XTAL\_CLK 是否携带毛刺，并在检测到毛刺后，产生数字系统复位信号，复位包括 RTC 在内的整个数字电路。

\section{功能描述}

\subsection{时钟毛刺检测}

ESP32-C3 的毛刺检测模块将对输入芯片的 XTAL\_CLK 时钟信号进行检测，当时钟的脉宽（a 或 b）小于 3 ns 时，将认为检测到毛刺，触发毛刺检测信号，屏蔽输入的 XTAL\_CLK 时钟信号。

\begin{figure}[H]
    \centering
    \includegraphics[width=0.5\textwidth]{\modulefiles/figures/glitch-clock-width}
    \caption{XTAL\_CLK 脉宽}
    \label{fig:glitch-pulse-width}
\end{figure}

\subsection{复位}
当时钟毛刺检测电路检测到 XTAL\_CLK 上有影响电路正常工作的毛刺之后，如果 \hyperref[fielddesc:RTCCNTLGLITCHRSTEN]{RTC\_CNTL\_GLITCH\_RST\_EN} 位为1，将触发系统级复位。该位默认为开启复位状态。

\end{document}
